\chapter{The Deterministic Harness — תשתית דטרמיניסטית סביב הסקת המודל}
\section{מהו Harness ולמה הוא דטרמיניסטי}
בלב הארכיטקטורה שלנו עומד רעיון פשוט אך מכריע: את הסקת המודל יש לעטוף בשכבת תוכנה דטרמיניסטית הקרויה Harness. מודל שפה הוא רכיב הסתברותי מטבעו — בהינתן אותו prompt הוא עשוי להחזיר פלטים שונים, להמציא תחביר או לחרוג מהמבנה הצפוי. לכן אנו מפרידים בחדות בין שתי שכבות אחריות. ה-Agents של CrewAI — Planner, Research, Writer, Reviewer ו-LaTeX — מבצעים את העבודה האינטלקטואלית הדורשת judgment: תכנון, סינתזה של מחקר, כתיבת פרוזה ועריכה. ה-Harness, לעומתם, מבצע את העבודה השבירה והניתנת לאימות באמצעות קוד Python קבוע: reconciliation של ציטוטים, יצירת graph, validation, escaping של תווים ו-compilation של ה-PDF. ההשראה לכך מגיעה מתפיסת ה-production harness, שבה Harness אמיתי עוטף את קריאות המודל ב-tools, ב-context, ב-parsing, ב-validation, ב-observability וב-guardrails. ההיגיון הוא שכל פעולה שתוצאתה חייבת להיות נכונה ובת-שחזור אסור שתישען על אִלתור של מודל. כך פלט ה-Agent נשמר כ-artifact מובנה ובר-בדיקה — למשל book\_plan או manuscript — וה-Harness צורך אותו, מאמת אותו וממיר אותו לתוצר סופי. תכן זה ממקם את אי-הוודאות במקום אחד ומבוקר: בתוך הסקת המודל בלבד. כל מה שמסביב הופך לצינור הנדסי צפוי, שניתן להריץ שוב ושוב ולקבל בכל פעם בדיוק אותה התנהגות, וכך מערכת יצירתית במהותה הופכת למערכת אמינה תפעולית.

\section{ציטוטים, graph ו-validation כרכיבים דטרמיניסטיים}
שלוש משימות בולטות מודגמות במערכת כרכיבים דטרמיניסטיים מובהקים — ולא כ-Agents. הראשונה היא ניהול הציטוטים. אסור שמודל ימציא רשומות bibliography או citation keys, ולכן ה-Citation Manager הוא קוד Python: הוא קורא source registry אצור ומבוקר, מייצר קובץ references.bib תקין תחבירית עם escaping בטוח, ומבצע reconciliation המוודא שכל citation marker ב-manuscript נפתר למקור מוכר. אם מופיע מפתח שאינו מזוהה, התהליך נכשל בקול רם עוד לפני ה-compilation, ועיקרון מפתח כאן הוא idempotency: אותו registry מפיק תמיד בדיוק אותו .bib. המשימה השנייה היא יצירת ה-graph. השרטוט המתאר את ה-pipeline — Planner אל Research אל Writer אל Reviewer אל LaTeX ואל רכיבי הבנייה הדטרמיניסטיים — נוצר בקוד Python קבוע ולא על ידי Agent, שכן graph דטרמיניסטי הוא בר-שחזור וקל ל-validation. המשימה השלישית היא ה-validation עצמה. ה-Validator הדטרמיניסטי מוודא שכל הרכיבים הנדרשים קיימים לפני ה-compilation: cover metadata, TOC, chapters ו-sections, לפחות תמונה אחת, קובץ ה-graph, table, נוסחת display שאינה טקסט רגיל, מקטע BiDi עברית-אנגלית, citation keys תקפים וקובצי ה-LaTeX source. חלוקה זו עקבית עם תפיסת הקורס: עבודה השופטת איכות ומשמעות נשארת אצל Agent כמו ה-Reviewer, אך הבדיקה הטכנית — האם הכול מחובר, מתקמפל ותקין — חייבת להיות דטרמיניסטית, שכן זו בדיקה הנדסית הניתנת ל-testing אובייקטיבי ולא שיפוט תוכן.

\section{דטרמיניזם, reproducibility ופלט אמין}
הדטרמיניזם אינו קישוט הנדסי אלא תנאי יסוד לשלושה ערכים שבלעדיהם המערכת אינה ראויה לאמון. הראשון הוא reproducibility. כאשר רכיבי ה-Harness מתנהגים באופן צפוי, הרצה חוזרת על אותם artifacts מפיקה בדיוק את אותו פלט — אותו references.bib, אותו graph ואותו validation report. תכונה זו חיונית במיוחד בתוצר אקדמי מדורג, שבו בודק חייב להיות מסוגל לשחזר את הבנייה ולקבל את אותו final.pdf. הערך השני הוא testing. רכיב הסתברותי קשה לבדיקה, שכן אין לו פלט יחיד נכון; רכיב דטרמיניסטי, לעומתו, ניתן לכסות ב-unit tests מדויקים: לוודא ש-citation key שאינו מזוהה מדווח כ-unresolved, שתווים מיוחדים עוברים escaping, ושאותו registry מפיק שוב את אותו .bib. כך ה-quality gates — schema validation, citation validation, asset validation ו-LaTeX validation — הופכים לבדיקות אוטומטיות אמינות במקום הסתמכות על שיפוט אנושי. הערך השלישי הוא trustworthy output. בכך שאנו ממקמים את אי-הוודאות אך ורק בתוך הסקת המודל ועוטפים אותה ב-validation וב-guardrails, אנו מבטיחים שתקלה של Agent — מפתח שגוי, נכס חסר או נוסחה פגומה — נתפסת מוקדם ובאופן גלוי, ואינה מתפשטת בשקט אל המסמך הסופי. שילוב היצירתיות ההסתברותית של המודל עם תשתית דטרמיניסטית קפדנית הוא שמאפשר למערכת להיות בו-זמנית גמישה בתוכן ואמינה בתוצר, וזהו בדיוק האיזון שתזמור סוכנים מקצועי שואף אליו.

לקריאה נוספת בנושאי פרק זה, ראו \cite{crewai_docs}.

